% ============================================================
% Section 6: Integration with the BX3 Framework
% Recursive Spawning Protocol — v6
% ============================================================
\section{Integration with the BX3 Framework}
\label{sec:rs-integration}

The Recursive Spawning Protocol is not a module composed alongside the BX3 Framework; it is a structural instantiation of it. Every RSP mechanism maps onto exactly one BX3 layer, and each such mapping is minimal and structurally necessary. The three-layer architecture is not a design choice among plausible alternatives — it is the unique minimal sufficient substrate for RSP's correctness guarantees. This section specifies the three mappings and establishes their structural necessity.

% ----------------------------------------------------------
\subsection{Purpose Layer: Accountability Anchor}
\label{sec:rs-purpose-anchor}

The Purpose Layer performs two roles in RSP that no other layer can perform.

\medskip
\noindent\textbf{Exclusive origin of spawn authorization.} At root provisioning, the Purpose Layer defines the spawn authorization policy $P_{\mathrm{spawn}}$ and records it in the Fact Layer authorization ledger. $P_{\mathrm{spawn}}$ encodes two mandatory pre-conditions for any child spawn: (i) the child operating scope must be a strict descendant refinement of the parent scope ($R(n_{i+1}) \subset R(n_i)$), and (ii) the spawn must be operationally necessary — the unmet condition must be unsatisfiable within the parent's existing scope. The Bounds Engine evaluates proposed spawns against $P_{\mathrm{spawn}}$; the Fact Layer enforces it. Neither originates it. No machine actor may issue a spawn warrant in the absence of a matching Purpose Layer authorization record. The separation is architectural, not organizational.

\medskip
\noindent\textbf{Exclusive terminus of escalation.} When the chain reaches $D_{\mathrm{ESCALATE}}$, the protocol compulsory-escalates to the human anchor $h(n)$. The human anchor is non-collapsing: for all nodes $n_i$ in chain $C$, $h(n_i) = h(n_0)$. The human issues one of three directives — \texttt{ACK}, \texttt{RESOLVE}, or \texttt{REJECT} — and no machine actor may issue any of these. The chain terminates in human judgment, never in autonomous resolution. This is the structural expression of the upstream BX3 accountability guarantee: systems built on the BX3 Framework fail upward into human consciousness, never downward into autonomous chaos.

The structural necessity of the Purpose Layer follows from what its roles demand: the exercise of intent and the bearing of final accountability — capabilities structurally reserved for the human principal. The Bounds Engine cannot originate $P_{\mathrm{spawn}}$ because it operates within the operating range defined by the Purpose Layer; it has no mandate to define what counts as authorized purpose. The Fact Layer cannot serve as the escalation terminus because it is a deterministic enforcement layer, not a judgment layer. Removing the Purpose Layer collapses the exclusive human terminus and enables self-authorized spawn by machine actors.

% ----------------------------------------------------------
\subsection{Bounds Engine: Depth Enforcement}
\label{sec:rs-bounds-depth}

The Bounds Engine enforces depth constraints that prevent unbounded chain growth. Critically, it cannot exceed these constraints under any operational circumstances.

\medskip
\noindent\textbf{Depth evaluation at spawn.} At each spawn event, the Bounds Engine evaluates current chain depth $|C|$ against $D_{\mathrm{max}}$, anchored in the Fact Layer authorization ledger. If $|C| \geq D_{\mathrm{max}}$, the Bounds Engine does not execute the spawn. It transitions to \texttt{ESCALATING} state and initiates the Bailout Protocol, propagating the exception upward to $h(n)$. This enforcement is deterministic: the Bounds Engine has no discretion to exceed $D_{\mathrm{max}}$ regardless of urgency, operational exigency, or the importance of the unmet condition. The depth constraint is a hard architectural boundary, not a configurable parameter.

\medskip
\noindent\textbf{Scope refinement evaluation.} The Bounds Engine also evaluates the scope subset condition $R(n_{i+1}) \subset R(n_i)$ at each spawn event. This evaluation is a precondition for the Fact Layer pre-commit hook; it is not sufficient alone because the scope refinement condition is independently enforced by the pre-commit hook. This separation reflects the depth-limit insufficiency theorem established in Section~\ref{sec:rs-drift-problem}: depth limits alone cannot prevent scope creep drift, so both depth enforcement and scope refinement enforcement are required as structurally distinct operations.

\medskip
\noindent\textbf{Structural reinforcement by the Fact Layer.} The Bounds Engine's depth evaluation is reinforced by the Fact Layer pre-commit hook. Before any spawn operation is recorded, the Fact Layer independently verifies: (a) $R(n_{i+1}) \subseteq R(n_i)$ (scope subset), and (b) $\mathrm{depth}(n_{i+1}) \leq D_{\mathrm{max}}$ (depth bound). The Fact Layer rejects any spawn violating either condition. The pre-commit hook makes the depth bound and scope refinement constraint structural invariants — not policy conventions overridable in exigent circumstances.

The Bounds Engine is necessary because the Fact Layer can reject a spawn exceeding the depth bound, but without the Bounds Engine there is no layer to evaluate whether a spawn is operationally necessary within the authorized scope, nor any layer to manage the transition to \texttt{ESCALATING} state when the bound is reached. Removing the Bounds Engine eliminates the constrained execution evaluation layer.

% ----------------------------------------------------------
\subsection{Fact Layer: Forensic Ledger}
\label{sec:rs-fact-ledger}

The Fact Layer provides deterministic enforcement and forensic auditability. It maintains the append-only ledger that makes the RSP chain auditable at any future time and enforces the write protocol constraints that make RSP structurally tamper-evident.

\medskip
\noindent\textbf{Append-only forensic ledger.} The Fact Layer records every spawn event, state transition, Purpose Layer directive received, and unresolved exception. Each entry is timestamped, attributed to the generating node, and hash-chained to the preceding entry. Hash-chaining ensures no entry can be inserted, deleted, reordered, or altered after recording: any tampering breaks hash continuity and is detectable by any observer with ledger read access. The ledger is the single source of truth for all authorized chain state — the runtime artifact that distinguishes RSP chains from conventional delegation chains, which exist only as forensic reconstructions attempted after the fact.

\medskip
\noindent\textbf{Immutable parent pointer.} Once $\alpha(n)$ is established at node provisioning, no actor may modify it. The Fact Layer write protocol rejects any attempt to overwrite $\alpha(n)$. Chain topology changes occur only through authorized spawn operations recorded in the ledger; there is no mechanism for retroactive pointer manipulation. This eliminates the re-parenting drift pathway (Failure Mode 3) by making retroactive chain rewriting structurally impossible regardless of the requestor's identity or privilege level.

\medskip
\noindent\textbf{Pre-commit hook.} The pre-commit hook verifies conditions (a) and (b) before any spawn is recorded. Rejected spawns are logged; the Bounds Engine does not execute them; the protocol transitions to \texttt{ESCALATING} if the unmet condition persists. The pre-commit hook operates identically on all machine actors; there is no privileged actor capable of bypassing it. Even administrative compromise of a machine actor cannot produce a bypass, because the hook is enforced at the Fact Layer write protocol level.

The Fact Layer is structurally necessary because without the Fact Layer write protocol, immutability is a promise rather than a constraint. A mutable parent pointer enables all four drift failure modes. Without the pre-commit hook, scope refinement and depth bound are policy documentation — enforceable only until an actor with sufficient privilege decides they are inconvenient. Removing the Fact Layer eliminates both the immutable parent pointer and the forensic ledger.

% ----------------------------------------------------------
\subsection{Structural Minimality}
\label{sec:rs-minimality}

The three-layer mapping is minimal and jointly sufficient. Removing any layer causes at least one RSP correctness property to fail:

\begin{itemize}
  \item \textbf{Removing Purpose Layer:} The exclusive human terminus is eliminated; the chain can terminate in machine actors; spawn authorization becomes self-issued by machine actors; the human anchor is no longer a compulsory terminus.
  \item \textbf{Removing Bounds Engine:} The constrained execution evaluation layer is eliminated; no layer evaluates spawn necessity or manages the transition from pre-commit rejection to mandatory escalation.
  \item \textbf{Removing Fact Layer:} The immutable parent pointer is eliminated; retroactive chain manipulation becomes possible; scope refinement and depth bound degrade from structural invariants to overridable policy conventions.
\end{itemize}

No proper subset of $\{$\ Purpose Layer, Bounds Engine, Fact Layer$\ \}$ achieves the same guarantees. This is not an empirical observation but a structural consequence of the layers' roles: the Purpose Layer's unique access to intent and final accountability, the Bounds Engine's unique access to constrained execution evaluation, and the Fact Layer's unique access to deterministic write enforcement and tamper-evident storage. Each layer's necessity is established by what no other layer can perform.

% ============================================================
% Section 7: Correctness Properties
% ============================================================
\section{Correctness Properties}
\label{sec:rs-correctness}

This section establishes three correctness properties for the Recursive Spawning Protocol: drift prevention, chain integrity, and termination. Each is stated formally, proved sketchedly, and connected to the BX3 Framework layer that enforces it. Together these properties constitute the RSP's overall guarantee: a spawn chain is an accountable, bounded, and auditable refinement process that terminates in human judgment.

% ----------------------------------------------------------
\subsection{Drift Prevention}
\label{sec:rs-drift-prevention}

\medskip
\noindent\textbf{Theorem (Drift Prevention).}
Let $C = \langle n_0, n_1, \ldots, n_k \rangle$ be a Recursive Spawning chain rooted at $n_0$ with human anchor $h(n_0) = h(n_k)$. Then for all $i$ ($0 \leq i \leq k$), the operating scope $R(n_i)$ is a descendant refinement of the intent of $h(n_0)$. Equivalently: no node in $C$ can exhibit behavior outside the authorized intent of its human anchor.

\medskip
\noindent\textbf{Proof sketch.} The proof proceeds by structural induction on chain depth, where each inductive step is enforced at runtime by the conjunction of three BX3-layer guarantees.

\medskip
\textit{Base case ($i=0$).}
$R(n_0)$ is defined by $h(n_0)$ at provisioning. By the Purpose Layer's definition of authorized operating scope, $R(n_0)$ is a descendant refinement of $h(n_0)$'s intent by construction. No drift exists at the root.

\medskip
\textit{Inductive step.}
Assume $R(n_i)$ is a descendant refinement of $h(n_0)$'s intent for some $i$ with $0 \leq i < k$. For $n_i$ to spawn $n_{i+1}$, two conditions must hold before the spawn is recorded:

\begin{enumerate}
  \item[(a)] $R(n_{i+1}) \subseteq R(n_i)$ — the scope refinement constraint, required by the Purpose Layer's spawn authorization policy $P_{\mathrm{spawn}}$.
  \item[(b)] $\mathrm{depth}(n_{i+1}) \leq D_{\mathrm{max}}$ — the depth constraint, enforced by the Bounds Engine and independently verified by the Fact Layer pre-commit hook.
\end{enumerate}

The Fact Layer pre-commit hook rejects any spawn violating (a) or (b). No actor, including the Bounds Engine, can execute a rejected spawn. Therefore $R(n_{i+1}) \subseteq R(n_i)$ holds by structural enforcement.

By the inductive hypothesis, $R(n_i)$ is a descendant refinement of $h(n_0)$'s intent. Since $R(n_{i+1}) \subseteq R(n_i)$, transitivity of subset gives that $R(n_{i+1})$ is also a descendant refinement of $h(n_0)$'s intent. Hence $n_{i+1}$ does not drift. By induction, no node in $C$ drifts. $\square$

\medskip
The drift prevention guarantee rests on the \emph{joint necessity} of three structural conditions. No single condition suffices; all three are required simultaneously.

\medskip
\noindent\textbf{Condition 1 — Immutable parent pointer (Fact Layer).}
$\alpha(n)$ is established once at provisioning and is thereafter immutable under the Fact Layer write protocol. If $\alpha(n)$ were mutable, a drifted actor could rewrite its ancestry to launder authorization retroactively. The immutable parent pointer is the structural constraint that makes the chain a verifiably honest record of the delegation lineage, not a post-hoc constructed narrative.

\medskip
\noindent\textbf{Condition 2 — Forensic ledger (Fact Layer).}
Every scope assignment is recorded in the append-only, hash-chained ledger. The ledger makes the inductive proof empirically verifiable: each step of the refinement chain is a recorded, tamper-evident ledger entry. Without the forensic ledger, the chain's scope history is a memory trace subject to the same failure modes as the mutable parent pointer it was designed to correct.

\medskip
\noindent\textbf{Condition 3 — Chain terminates at Purpose Layer (Purpose Layer).}
The exclusive human terminus closes the inductive argument at a finite bound. There is no machine-actor alternative terminus that could absorb exceptions and allow drift to compound silently. The human anchor $h(n)$ is non-collapsing: $h(n_i) = h(n_0)$ for all $n_i \in C$. If the chain could terminate in a machine actor, an autonomous exception absorber could issue a terminal directive and the chain would terminate in autonomous resolution — the upstream BX3 accountability guarantee would be violated.

\medskip
The three conditions are jointly necessary and jointly sufficient. If any one is absent, drift prevention fails:
\begin{itemize}
  \item Without the immutable parent pointer (Condition 1), retroactive scope laundering becomes possible.
  \item Without the forensic ledger (Condition 2), scope refinement history cannot be independently verified.
  \item Without the Purpose Layer terminus (Condition 3), the chain can terminate in autonomous resolution, making drift unobservable.
\end{itemize}
With all three present, drift is architecturally impossible: the immutable parent pointer prevents retroactive scope manipulation, the forensic ledger makes every refinement step auditable, and the Purpose Layer terminates the chain in human judgment.

% ----------------------------------------------------------
\subsection{Chain Integrity}
\label{sec:rs-chain-integrity}

\medskip
\noindent\textbf{Theorem (Chain Integrity).}
For any Recursive Spawning chain $C = \langle n_0, n_1, \ldots, n_k \rangle$, the parent pointer sequence $\alpha(n_i) = n_{i-1}$ holds for all $1 \leq i \leq k$, and no node configuration in $C$ has been modified after provisioning.

\medskip
\noindent\textbf{Proof sketch.} The proof has two components: topological integrity and node immutability.

\medskip
\textit{Topological integrity.}
Each spawn event $n_{i-1} \xrightarrow{\mathrm{spawn}} n_i$ establishes $\alpha(n_i) = n_{i-1}$ as an immutable Fact Layer property at provisioning time. The Fact Layer write protocol rejects any operation attempting to modify $\alpha(n_i)$ after provisioning. Since the chain is constructed exclusively through authorized spawn events and no parent pointer is modified after establishment, $\alpha(n_i) = n_{i-1}$ holds for all $i \geq 1$ by construction.

The parent pointer is not merely access-controlled — it is structurally immutable. There is no privileged actor, including the Purpose Layer after provisioning, capable of overwriting $\alpha(n)$. This eliminates the re-parenting drift pathway (Failure Mode 3) by making retroactive chain rewriting structurally impossible. Chain topology changes occur only through authorized spawn events recorded in the ledger; there is no mechanism for retroactive pointer manipulation.

\medskip
\textit{Node immutability.}
Each node $n_i$ is provisioned with a fixed configuration $(R(n_i), P_{\mathrm{spawn}}(n_i), h(n_i))$ established by the Purpose Layer and recorded in the Fact Layer authorization ledger. The Fact Layer enforces immutability of these parameters after provisioning. The write protocol rejects any modification request from any actor — including the Bounds Engine and the Purpose Layer after provisioning. A \texttt{RESOLVE} directive redirects node behavior but is recorded as a new ledger entry (a state transition), not as an alteration of the provisioned configuration.

The forensic ledger's hash-chaining provides evidentially for chain integrity: any attempt to insert a falsified spawn event or reorder existing entries breaks hash continuity and is detectable by any observer with ledger read access. Chain integrity is a property of the entire historical record, enforceable at any future time. $\square$

% ----------------------------------------------------------
\subsection{Termination}
\label{sec:rs-termination}

\medskip
\noindent\textbf{Theorem (Termination).}
Every Recursive Spawning chain $C = \langle n_0, n_1, \ldots, n_k \rangle$ terminates — either at a resolution state or at the escalation threshold — within a finite number of spawn steps bounded by $\max(D_{\mathrm{max}}, D_{\mathrm{ESCALATE}})$.

\medskip
\noindent\textbf{Proof sketch.} The proof establishes two lemmas and combines them.

\medskip
\textit{Lemma 1 (Depth-bounded growth).}
The chain grows only through spawn events. By the Bounds Engine depth enforcement and the Fact Layer pre-commit hook, any spawn event producing $|C| \geq D_{\mathrm{max}}$ is rejected. Therefore the chain cannot exceed depth $D_{\mathrm{max}} - 1$ through autonomous spawn. Rejected spawns are logged; the Bounds Engine transitions to \texttt{ESCALATING}. Hence autonomous chain growth is bounded by $D_{\mathrm{max}}$.

\medskip
\textit{Lemma 2 (Escalation-triggered halt).}
If the chain reaches $D_{\mathrm{ESCALATE}}$ before resolving the exception condition, the protocol enters \texttt{ESCALATING} state and suspends all autonomous spawn activity. The Bounds Engine awaits a directive from $h(n_k)$ — the human anchor. No further spawn events occur until $h(n_k)$ issues \texttt{ACK}, \texttt{RESOLVE}, or \texttt{REJECT}. Since $h(n_k) = h(n_0)$ is a human principal, the directive arrives in finite time. The \texttt{REJECT} directive terminates the chain definitively; \texttt{RESOLVE} redirects behavior and records a new ledger entry; \texttt{ACK} permits continued operation under continued Purpose Layer authorization.

\medskip
\textit{Theorem conclusion.}
Combining Lemma 1 and Lemma 2: the chain grows autonomously at most $D_{\mathrm{max}} - 1$ spawn steps. In all execution paths, the chain reaches a terminal state within $\max(D_{\mathrm{max}}, D_{\mathrm{ESCALATE}})$ spawn steps. Since both $D_{\mathrm{max}}$ and $D_{\mathrm{ESCALATE}}$ are finite integers defined at provisioning, the chain terminates in finite time. $\square$

The termination guarantee also precludes infinite spawn-escalate cycles. When $D_{\mathrm{max}}$ is reached, the Fact Layer blocks further spawn and triggers mandatory escalation. The human anchor's \texttt{REJECT} directive terminates the chain. There is no mechanism by which the chain can circumvent this bound and re-enter autonomous spawn — the Fact Layer pre-commit hook is structural, not configurable at runtime by any machine actor.

% ----------------------------------------------------------
\subsection{Collective Guarantee}
\label{sec:rs-collective-guarantee}

Together, these three properties guarantee that a Recursive Spawning chain is an accountable, bounded, and auditable refinement process:

\begin{itemize}
  \item \textbf{Drift prevention} ensures scope refinement is the only authorized direction of chain growth. Each child scope is a strict subset of its parent, traceable to the human anchor's intent through the forensic ledger. Without drift prevention, successive scope expansions could gradually migrate the chain's operating scope away from the human anchor's intent — an autonomy runway within an otherwise bounded-depth chain.
  \item \textbf{Chain integrity} ensures the chain topology is stable and tamper-evident. The parent pointer sequence is immutable, node configurations are unmodifiable after provisioning, and the forensic ledger's hash-chaining makes any tampering detectable. Without chain integrity, an adversarial actor could manipulate chain parentage or configuration to redirect accountability or obscure the provenance of a spawned node — a form of accountability laundering.
  \item \textbf{Termination} ensures the chain cannot grow indefinitely and cannot sustain an infinite spawn-escalate cycle. The depth bound is enforced structurally by the Fact Layer pre-commit hook, and the escalation path terminates in human judgment, never in autonomous resolution. Without termination, an unresolvable exception condition could drive unbounded chain growth exceeding any human's supervisory capacity.
\end{itemize}

Together, Theorems~\ref{sec:rs-drift-prevention}, \ref{sec:rs-chain-integrity}, and~\ref{sec:rs-termination} establish that a Recursive Spawning chain governed by the BX3 Framework is a provably accountable, integrity-preserving, and finite lineage of agents — suitable for deployment in high-stakes operational contexts where auditability and correctness are non-negotiable requirements.
